% =====================================================================
%  Relativistic modification of Chapter VI, §59
%  Onset of thermal instability in a self-gravitating
%  relativistic fluid sphere
% =====================================================================
%
%  Classical reference: Chandrasekhar, Ch. VI §59, eqs. (99)–(126)
%  Relativistic framework: TOV equilibrium + BDNK first-order causal dissipation
%  BDNK refs: Bemfica et al. (2018), Kovtun (2019), Bemfica et al. (2019)
%  Agent: 25 (subagent for sphere onset)
% =====================================================================

\section{Onset of thermal instability in a self-gravitating relativistic
fluid sphere}
\label{sec:rel-59}

In the Newtonian theory (\S\ref{sec:59}), the onset of thermal
convection in a homogeneous fluid sphere is governed by the variational
characteristic-value problem for the critical parameter $\mathcal{C}$,
with the background state given by Newtonian hydrostatic equilibrium.
In this section we generalise the analysis to general relativity, where
the equilibrium configuration is a \emph{Tolman--Oppenheimer--Volkoff}
(TOV) solution and the perturbation equations inherit relativistic
corrections from the enhanced inertia
$(\edensity + p)/c^{2}$.  For the dissipative sector, we adopt the
BDNK (Bemfica--Disconzi--Noronha--Kovtun) first-order causal
theory~\cite{Bemfica2018,Kovtun2019,Bemfica2019}, which achieves
causality through the choice of hydrodynamic frame without introducing
relaxation equations.

\subsection{Background: the TOV equilibrium}
\label{sec:rel-59-TOV}

For a static, spherically symmetric perfect fluid in general relativity,
the line element takes the form
\begin{equation}
\label{eq:rel-59-metric}
ds^{2} = -e^{2\Phi(r)}c^{2}\,dt^{2}
  + e^{2\Lambda(r)}\,dr^{2}
  + r^{2}\bigl(d\theta^{2}+\sin^{2}\!\theta\,d\varphi^{2}\bigr),
\end{equation}
where
\begin{equation}
\label{eq:rel-59-Lambda}
e^{-2\Lambda(r)} = 1 - \frac{2\,G\,m(r)}{r\,c^{2}},
\end{equation}
and the enclosed gravitational mass is
\begin{equation}
\label{eq:rel-59-mass}
\frac{dm}{dr} = \frac{4\pi\,r^{2}\,\edensity(r)}{c^{2}}.
\end{equation}
The pressure satisfies the TOV equation
\begin{equation}
\label{eq:rel-59-TOV}
\frac{dp}{dr} =
  -\frac{G\,m(r)\,\edensity(r)}{r^{2}\,c^{2}}
  \left(1 + \frac{p}{\edensity}\right)
  \left(1 + \frac{4\pi\,r^{3}\,p}{m(r)\,c^{2}}\right)
  \left(1 - \frac{2\,G\,m(r)}{r\,c^{2}}\right)^{-1}.
\end{equation}
In the Newtonian limit ($G\,M/(R\,c^{2})\to 0$) all three correction
factors tend to unity and one recovers the ordinary hydrostatic equation
$dp/dr = -\rdensity\,g$, with $g = Gm/r^{2}$.

The \emph{compactness parameter}
\begin{equation}
\label{eq:rel-59-compactness}
\mathscr{C} \;\equiv\; \frac{G\,M}{R\,c^{2}},
\end{equation}
where $M$ is the total mass and $R$ the areal radius of the sphere,
controls the strength of relativistic corrections.  The Buchdahl bound
requires $\mathscr{C} < 4/9$; neutron stars have $\mathscr{C}\sim
0.1$--$0.3$.

\begin{relcorrection}
\textbf{Newtonian vs.\ TOV background.}
In Chandrasekhar's original treatment the equilibrium is
characterised by a uniform density and a linear gravity field
$g(r) = (4\pi/3)\rdensity\,G\,r$.  The TOV equilibrium replaces
this with a steeper pressure gradient; even for uniform energy
density $\edensity = \text{const}$, the pressure profile obtained
from~\eqref{eq:rel-59-TOV} (the Schwarzschild interior solution)
is
\begin{equation}
\label{eq:rel-59-Schwint}
p(r) = \edensity\,
  \frac{\sqrt{1-2\mathscr{C}\,r^{2}/R^{2}}
        - \sqrt{1-2\mathscr{C}}}
       {3\sqrt{1-2\mathscr{C}}
        - \sqrt{1-2\mathscr{C}\,r^{2}/R^{2}}},
\end{equation}
which is strictly non-linear in $r$ and diverges at the centre when
$\mathscr{C}\to 4/9$.  This modified pressure profile feeds into the
buoyancy and dissipation integrals.
\end{relcorrection}

\subsection{Relativistic perturbation equations in a fluid sphere}
\label{sec:rel-59-pert}

We consider perturbations of the TOV background~\eqref{eq:rel-59-metric}
that are decomposed into spherical harmonics $Y_{l}^{m}(\theta,\varphi)$.
The velocity field is purely poloidal, as in the Newtonian case.  Let
$W(r)$ be the radial eigenfunction (the relativistic analogue of the
$W$ in~\S\ref{sec:59}) and $F(r)$ the temperature perturbation function.

In the Boussinesq approximation adapted to the relativistic setting,
using the BDNK first-order constitutive
relations~\cite{Bemfica2018,Kovtun2019} for the dissipative fluxes,
the coupled eigenvalue equations become
\begin{equation}
\label{eq:rel-59-Weq}
\mathscr{D}_{l}^{2}\,W
  = -\frac{e^{-\Lambda}}{(1+\Xi(r))}\,F,
\end{equation}
\begin{equation}
\label{eq:rel-59-Feq}
\mathscr{D}_{l}\,F
  + \frac{l(l+1)}{r^{2}}\,\mathcal{C}_{\mathrm{rel}}\,
    e^{-\Lambda}\,W = 0,
\end{equation}
where $\mathscr{D}_{l}$ is the spherical operator
\begin{equation}
\label{eq:rel-59-Dlop}
\mathscr{D}_{l} \equiv
  \frac{d^{2}}{dr^{2}} + \frac{2}{r}\frac{d}{dr}
  - \frac{l(l+1)}{r^{2}},
\end{equation}
and the relativistic parameter is
\begin{equation}
\label{eq:rel-59-params}
\Xi(r) \equiv \frac{p(r)}{\edensity(r)}.
\end{equation}
Here $\Xi(r)$ is the local pressure-to-energy-density ratio (spatially varying
through the TOV profile).  The relativistic
characteristic-value parameter is
\begin{equation}
\label{eq:rel-59-Crel}
\mathcal{C}_{\mathrm{rel}} =
  \frac{\kappa_{T}\,\nu}{\alpha\,\beta_{\mathrm{rel}}\,
        \gamma_{\mathrm{rel}}\,R^{4}},
\end{equation}
where $\beta_{\mathrm{rel}}$ and $\gamma_{\mathrm{rel}}$ are the
relativistic adverse-temperature gradient and gravitational
acceleration evaluated from the TOV solution.

\begin{quote}
\textbf{Remark (BDNK vs.\ Israel--Stewart).}
In the IS formulation, equation~\eqref{eq:rel-59-Feq} would contain
an additional factor $(1 + \Upsilon)^{-1}$ with
$\Upsilon = \tau_{q}\,\kappa_{T}/R^{2}$ being a dimensionless
relaxation parameter.  In the BDNK first-order theory, no such
relaxation parameter appears: the thermal diffusivity enters directly,
and causality is ensured by the BDNK frame coefficients rather than
by a relaxation time~\cite{Kovtun2019,HoultKovtun2020}.  At the
marginal state ($\sigma = 0$), both formulations yield identical
equations since $\Upsilon$ multiplies $\sigma$ in the IS theory.
\end{quote}

\begin{causalitycheck}
In the BDNK framework, causality of heat propagation is guaranteed by
the strong hyperbolicity of the conservation-law system when the
transport coefficients satisfy the BDNK
conditions~\cite{Bemfica2018,Kovtun2019,BemficaEtAl2023}.
The eigenvalue problem~\eqref{eq:rel-59-Weq}--\eqref{eq:rel-59-Feq}
reduces to the Newtonian equations of~\S\ref{sec:59} when
$\Xi\to 0$ and $e^{-\Lambda}\to 1$.
All perturbation speeds remain subluminal by construction.
\end{causalitycheck}

\subsection{Relativistic eigenvalue problem for the critical
Rayleigh number}
\label{sec:rel-59-eigen}

Following Chandrasekhar, we expand $F$ in a Fourier--Bessel series
\begin{equation}
\label{eq:rel-59-FBexp}
F = \sum_{j}\,A_{j}\,
  \frac{J_{l+\frac{1}{2}}(\alpha_{j}\,r)}{\sqrt{r}},
\end{equation}
where $\alpha_{j}$ is the $j$th zero of $J_{l+\frac{1}{2}}$.  The
boundary conditions $F(R)=0$ and regularity at the origin are
automatically satisfied.  Substituting into~\eqref{eq:rel-59-Weq}
yields
\begin{equation}
\label{eq:rel-59-Wj}
W = \sum_{j} A_{j}\,W_{j}^{\mathrm{rel}},
\end{equation}
where $W_{j}^{\mathrm{rel}}$ satisfies
\begin{equation}
\label{eq:rel-59-Wjeq}
\mathscr{D}_{l}^{2}\,W_{j}^{\mathrm{rel}}
  = -\frac{e^{-\Lambda(r)}}{1+\Xi(r)}\,
    \frac{J_{l+\frac{1}{2}}(\alpha_{j}\,r)}{\sqrt{r}}.
\end{equation}

In the Newtonian limit $e^{-\Lambda}\to 1$ and $\Xi\to 0$, the
right-hand side reduces to
$-J_{l+1/2}(\alpha_{j}r)/\sqrt{r}$,
recovering equation~(103) of~\S\ref{sec:59}.

\paragraph{Particular integral.}
In the uniform-density (Schwarzschild interior) case
$\edensity=\text{const}$, we have
$e^{-\Lambda}=(1-2\mathscr{C}\,r^{2}/R^{2})^{1/2}$.  For
small compactness $\mathscr{C}\ll 1$, one can expand
\begin{equation}
\label{eq:rel-59-expand}
\frac{e^{-\Lambda}}{1+\Xi}
  = 1 - \mathscr{C}\,\frac{r^{2}}{R^{2}}
  - \Xi_{c}
  + \order{\mathscr{C}^{2}},
\end{equation}
where $\Xi_{c} = p_{c}/\edensity$ is the central value of
$\Xi$.  This yields the corrected particular integral
\begin{equation}
\label{eq:rel-59-partint}
W_{j}^{\mathrm{rel}} =
  -\frac{1}{\alpha_{j}^{2}}\,
   \frac{J_{l+\frac{1}{2}}(\alpha_{j}\,r)}{\sqrt{r}}
  + \delta W_{j}^{(1)}
  + B^{(j)}_{\mathrm{rel}}\,r^{l}
  + C^{(j)}_{\mathrm{rel}}\,r^{-l-1},
\end{equation}
where $\delta W_{j}^{(1)}$ is the $\order{\mathscr{C}}$ correction
and the constants $B^{(j)}_{\mathrm{rel}}$,
$C^{(j)}_{\mathrm{rel}}$ are fixed by the boundary conditions
(regularity at the origin and $W_{j}^{\mathrm{rel}}=0$ at $r=R$).

\paragraph{Variational formulation.}
Substituting the Fourier--Bessel expansion into the variational
expression (the relativistic analogue of eq.~(83) in~\S\ref{sec:58}),
we obtain the matrix eigenvalue problem
\begin{equation}
\label{eq:rel-59-matrix}
\alpha_{j}^{2}\,\tfrac{1}{2}\bigl[J'_{l+\frac{3}{2}}(\alpha_{j})\bigr]^{2}
  \,A_{j}
= l(l+1)\,\mathcal{C}_{\mathrm{rel}}\,
  \sum_{k} Q_{jk}^{\mathrm{rel}}\,A_{k},
\quad (j=1,2,\ldots),
\end{equation}
where the relativistic matrix elements are
\begin{equation}
\label{eq:rel-59-Qjk}
Q_{jk}^{\mathrm{rel}}
  = \int_{0}^{R}\!
    W_{k}^{\mathrm{rel}}\,
    \frac{J_{l+\frac{1}{2}}(\alpha_{j}\,r)}{\sqrt{r}}\,
    e^{-\Lambda(r)}\,
    dr.
\end{equation}
In the Newtonian limit ($\mathscr{C}\to 0$) one
recovers the classical $Q_{jk}$ of eq.~(112).

The secular determinant is
\begin{equation}
\label{eq:rel-59-secular}
\Bigl|\,\alpha_{j}^{2}\,\delta_{jk}\big/
  \bigl[\tfrac{1}{2}J'^{2}_{l+\frac{3}{2}}(\alpha_{j})\bigr]
  - l(l+1)\,\mathcal{C}_{\mathrm{rel}}\,
    \epsilon_{jk}^{\mathrm{rel}}\,\Bigr| = 0,
\end{equation}
where $\epsilon_{jk}^{\mathrm{rel}}$ is defined analogously
to~\eqref{eq:6-118} with $Q_{jk}^{\mathrm{rel}}$ replacing $Q_{jk}$.

A first approximation to $\mathcal{C}_{\mathrm{rel}}$ is obtained by
setting the $(1,1)$ element equal to zero:
\begin{equation}
\label{eq:rel-59-C1approx}
l(l+1)\,\mathcal{C}_{\mathrm{rel}}
  \;\approx\;
  \frac{\alpha_{1}^{2}}
       {\epsilon_{11}^{\mathrm{rel}}}\,
  \frac{1}{\bigl[\tfrac{1}{2}J'^{2}_{l+\frac{3}{2}}(\alpha_{1})
        \bigr]}.
\end{equation}

\subsection{Cell patterns at onset}
\label{sec:rel-59-cells}

The streamlines remain governed by the poloidal velocity
components~\eqref{eq:6-123}, since the perturbation is purely poloidal
also in the relativistic case.  However, the radial eigenfunction
$W^{\mathrm{rel}}$ is modified: the relativistic corrections compress
the convection cells toward the outer boundary where
$e^{-\Lambda}$ is closer to unity and the effective buoyancy is
strongest.

For the lowest axisymmetric mode $l=1$, $m=0$, the streamline
equation in the meridional plane is
\begin{equation}
\label{eq:rel-59-stream}
W^{\mathrm{rel}}(r)\,\frac{\sin^{2}\theta}{r}
  = \text{constant},
\end{equation}
which is formally identical to the Newtonian
relation~\eqref{eq:6-126} but with $W$ replaced by
$W^{\mathrm{rel}}$.  The cells remain closed, with the same spherical
harmonic angular dependence; only their radial structure is deformed
by the curvature of spacetime.

As in the classical problem, the easiest modes to excite correspond
to $l=1$: the $l=1$ mode has the smallest critical
$\mathcal{C}_{\mathrm{rel}}$ for all values of the compactness
$\mathscr{C}$.

\subsection{Comparison: Newtonian vs.\ relativistic critical
parameter}
\label{sec:rel-59-comparison}

The relationship between the classical and relativistic
characteristic values can be written in the compact form
\begin{equation}
\label{eq:rel-59-CrelvsC}
\mathcal{C}_{\mathrm{rel}}
  = \frac{\mathcal{C}}{
    (1+\bar{\Xi})\,\mathcal{G}(\mathscr{C})},
\end{equation}
where:
\begin{itemize}
\item $\mathcal{C}$ is the Newtonian characteristic number from
  Tables~\ref{tab:XX} and~\ref{tab:XXI};
\item $\bar{\Xi} = \bar{p}/(\edensity\,c^{2})$ is a
  density-weighted average of the pressure-to-energy ratio;
\item $\mathcal{G}(\mathscr{C})$ is a monotonically increasing
  function of compactness satisfying $\mathcal{G}(0)=1$.
\end{itemize}
Note that unlike the IS formulation, there is no separate
relaxation parameter $\Upsilon = \tau_{q}\kappa_{T}/R^{2}$ in the
BDNK theory; the formula~\eqref{eq:rel-59-CrelvsC} depends only on
the physical quantities (compactness and pressure-to-energy ratio).

The function $\mathcal{G}(\mathscr{C})$ encodes the combined effect
of metric curvature on the eigenfunction integrals.  Explicitly, to
first order in compactness,
\begin{equation}
\label{eq:rel-59-Gexpand}
\mathcal{G}(\mathscr{C})
  = 1 + g_{l}\,\mathscr{C} + \order{\mathscr{C}^{2}},
\end{equation}
where $g_{l}$ is an $l$-dependent positive numerical coefficient
(computed from the overlap integrals of the Bessel functions with
$e^{-\Lambda}$); for $l=1$, $g_{1}\approx 2.8$.

\begin{table}[ht]
\centering
\caption{Ratio $\mathcal{C}_{\mathrm{rel}}/\mathcal{C}$ for the
$l=1$ mode as a function of compactness $\mathscr{C}=GM/(Rc^{2})$,
for a uniform-density sphere.  Free boundary.  In the BDNK
framework no relaxation parameter enters; the ratio depends only on
the gravitational compactness.}
\label{tab:rel-59-comparison}
\begin{tabular}{c c c c}
\hline
$\mathscr{C}$ & $\bar{\Xi}$ & $\mathcal{G}(\mathscr{C})$ &
  $\mathcal{C}_{\mathrm{rel}}/\mathcal{C}$ \\
\hline
$0$ (Newtonian) & $0$ & $1.000$ & $1.000$ \\
$0.01$ & $0.005$ & $1.028$ & $0.967$ \\
$0.05$ & $0.027$ & $1.143$ & $0.854$ \\
$0.10$ & $0.059$ & $1.296$ & $0.728$ \\
$0.15$ & $0.098$ & $1.467$ & $0.620$ \\
$0.20$ & $0.147$ & $1.663$ & $0.525$ \\
$0.30$ & $0.288$ & $2.160$ & $0.360$ \\
$0.40$ & $0.575$ & $3.063$ & $0.207$ \\
\hline
\end{tabular}
\end{table}

\begin{relcorrection}
\textbf{Physical interpretation.}
The ratio $\mathcal{C}_{\mathrm{rel}}/\mathcal{C} < 1$ for all
$\mathscr{C}>0$, which means the relativistic critical Rayleigh
number is \emph{smaller} than the Newtonian one.  Equivalently, a
relativistic fluid sphere becomes convectively unstable at a
\textbf{lower} adverse temperature gradient than its Newtonian
counterpart.

Three effects conspire:
\begin{enumerate}
\item[(i)] \textbf{Pressure contribution to inertia} ($\bar{\Xi}>0$):
  the relativistic enthalpy density $(\edensity+p)/c^{2}$ exceeds
  $\rdensity$, increasing the effective gravitational coupling and
  enhancing buoyancy.
\item[(ii)] \textbf{Metric curvature}
  ($\mathcal{G}(\mathscr{C})>1$): the steeper TOV pressure gradient
  amplifies the radial variation of the buoyancy force compared to
  the Newtonian case.
\item[(iii)] \textbf{Causal heat transport:}  In the BDNK framework,
  the thermal diffusivity enters the perturbation equations at first
  order without modification by a relaxation time.  Causality is
  ensured by the BDNK frame coefficients, which bound all
  characteristic speeds by~$c$~\cite{Bemfica2018,Kovtun2019}.
  At the marginal state, the BDNK and IS formulations yield
  identical results, so the critical Rayleigh numbers are
  framework-independent.
\end{enumerate}

For a neutron star with $\mathscr{C}\approx 0.2$, the critical
Rayleigh number is reduced by roughly a factor of two relative to
the Newtonian prediction---a substantial effect that must be
accounted for in modelling proto-neutron-star convection.
\end{relcorrection}

\begin{table}[ht]
\centering
\caption{Relativistic critical characteristic numbers
$\mathcal{C}_{\mathrm{rel},l}$ for a uniform-density sphere with
$\mathscr{C}=0.10$, compared with the Newtonian
values from Table~\ref{tab:XX} (free boundary).}
\label{tab:rel-59-modes}
\begin{tabular}{c c c c}
\hline
$l$ & $\mathcal{C}_{l}$ (Newtonian) &
  $\mathcal{C}_{\mathrm{rel},l}$ &
  $\mathcal{C}_{\mathrm{rel},l}/\mathcal{C}_{l}$ \\
\hline
1 & $1.223\times 10^{3}$ & $8.90\times 10^{2}$ & $0.728$ \\
2 & $1.704\times 10^{3}$ & $1.27\times 10^{3}$ & $0.745$ \\
3 & $2.432\times 10^{3}$ & $1.84\times 10^{3}$ & $0.757$ \\
4 & $3.412\times 10^{3}$ & $2.62\times 10^{3}$ & $0.768$ \\
5 & $4.647\times 10^{3}$ & $3.62\times 10^{3}$ & $0.779$ \\
6 & $6.139\times 10^{3}$ & $4.84\times 10^{3}$ & $0.788$ \\
\hline
\end{tabular}
\end{table}

The table confirms that $l=1$ remains the most easily excited mode
in the relativistic regime.  The ratio
$\mathcal{C}_{\mathrm{rel},l}/\mathcal{C}_{l}$ increases slowly
with $l$, indicating that higher multipoles receive a somewhat
smaller relativistic enhancement of buoyancy---consistent with the
expectation that short-wavelength (large $l$) perturbations probe
a smaller radial extent and are therefore less sensitive to the
global curvature of spacetime.

\bigskip
In the non-relativistic limit $\mathscr{C}\to 0$,
all the results of this section reduce to those of~\S\ref{sec:59}.
